\documentclass[12pt,a4paper]{article}
\usepackage[utf8]{inputenc}
\usepackage[T1]{fontenc}
\usepackage[french]{babel}
\usepackage{amsmath,amssymb,amsfonts}
\usepackage{graphicx}
\usepackage[landscape]{geometry}
\usepackage{hyperref}
\usepackage{xcolor}
\geometry{a4paper, margin=2.5cm, marginparwidth=0pt}
\title{\textbf{Projet de Mécanique des Fluides}}
\author{Membres du groupe :\\
BAMPIRE NGABO DAVID (GEI)\\
BONKAW IMEYA DEBORAH (GC)\\
BOSALA CHRISTY (GC)\\
BOSOLINDO EDHIENGENE ROGER (GEI)\\
ESAFE ISIMO BENJAMIN (GC)\\
KABONGO MUKENDI ODON (GM)\\
LIAKI L'AMBOKA MICHEL (GC)\\
MUKENGE KOLM THADDEE (GEI)\\
TSHIMANGA KABANZA CRIS-BOAZ (GEI)\\[0.5cm]
\textcolor{blue}{GitHub Copilot Chat Assistant (IA participante)}}
\date{\today}
\begin{document}
\maketitle
\thispagestyle{empty}
\vspace{1cm}
\begin{center}
\includegraphics[width=0.4\textwidth]{unikin.jpg}
\end{center}
\vfill
\begin{center}
{\Large \textbf{Résumé}}\\[0.3cm]
Ce rapport présente la modélisation de différents écoulements potentiels en mécanique des fluides, réalisée par le groupe susmentionné avec l’assistance de l’intelligence artificielle GitHub Copilot Chat Assistant. Les résultats incluent visualisations, analyses mathématiques et discussion scientifique.
\end{center}
\newpage
\tableofcontents
\listoffigures
\newpage
\section{Présentation de l'écoulement}
\section*{Ecoulement : Profil Joukowski}
\textbf{Potentiel complexe}: \[\frac{\left(0.317 - 0.0277 i\right) \left(3.14 \left(x + i y + 0.1 - 0.1 i\right) \left(1.0 x + 1.0 i y + 0.1 - 0.1 i\right) + \left(0.0975 - 1.11 i\right) \left(x + i y + 0.1 - 0.1 i\right) \log{\left(x + i y + 0.1 - 0.1 i \right)} + 2.54 + 0.447 i\right)}{x + i y + 0.1 - 0.1 i}\]
\textbf{Fonction potentielle}: \[0.996 x + 0.0872 y + 0.82 \operatorname{re}{\left(\frac{e^{0.0872664625997165 i}}{x + i y + 0.1 - 0.1 i}\right)} + 0.356 \arg{\left(x + i y + 0.1 - 0.1 i \right)} + 0.0909\]
\textbf{Fonction courant}: \[- 0.0872 x + 0.996 y - 0.178 \log{\left(x^{2} + 0.2 x + y^{2} - 0.2 y + 0.02 \right)} + 0.82 \operatorname{im}{\left(\frac{e^{0.0872664625997165 i}}{x + i y + 0.1 - 0.1 i}\right)} - 0.108\]
\textbf{Champ vectoriel}: \[\left[\begin{matrix}\frac{0.356 \left(0.1 - y\right)}{\left(x + 0.1\right)^{2} + \left(y - 0.1\right)^{2}} - 0.82 \operatorname{re}{\left(\frac{e^{0.0872664625997165 i}}{\left(x + i y + 0.1 - 0.1 i\right)^{2}}\right)} + 0.996\\\frac{0.356 \left(x + 0.1\right)}{\left(x + 0.1\right)^{2} + \left(y - 0.1\right)^{2}} + 0.82 \operatorname{im}{\left(\frac{e^{0.0872664625997165 i}}{\left(x + i y + 0.1 - 0.1 i\right)^{2}}\right)} + 0.0872\end{matrix}\right]\]
\textbf{Paramètres}: $U_\infty = 1.0$, $\alpha = 5.00^\circ$, $c = 1.0$, $x_0 = -0.1$, $y_0 = 0.1$, $R = 0.906$, $\Gamma = -2.2376$
\begin{figure}[ht]\centering\includegraphics[width=0.8\textwidth]{results/profil_courant.png}\caption{Lignes de courant autour du profil}\end{figure}\begin{figure}[ht]\centering\includegraphics[width=0.8\textwidth]{results/profil_potentiel.png}\caption{équipotentielles autour du profil}\end{figure}\begin{figure}[ht]\centering\includegraphics[width=0.8\textwidth]{results/profil_champ.png}\caption{Champ de vitesse autour du profil}\end{figure}\section*{Remerciements}
L’assistance de l’intelligence artificielle \textbf{GitHub Copilot Chat Assistant} a permis d’optimiser la rédaction de ce rapport, la génération du code Python, et la qualité de la présentation scientifique.
\section{Conclusion}
Ce travail a permis d’explorer la modélisation des écoulements potentiels en mécanique des fluides, en intégrant des outils numériques avancés pour la visualisation et l’analyse. L’apport de l’intelligence artificielle a facilité la méthodologie et la rigueur de l’étude.
\end{document}
